\subsection{Homogeneous Ensemble}
\label{meth:homo}

Constructing an ensemble from multiple ResNet variants represents a methodologically sound approach. 
ResNets of different depths, such as the 18-layer, 34-layer, and 50-layer ResNet, differ in model capacity, representational depth, and optimization behavior. 
Although all models are based on the same architectural principle, \citet{He_2016_CVPR} demonstrate that meaningful diversity exists within the ResNet family.

This diversity arises from the differing number of residual blocks. 
These blocks enable hierarchical refinement of feature representations. 
Deeper models, particularly the 50-layer ResNet, can therefore learn more complex and abstract representations. 
Shallower models such as the 18-layer ResNet, in contrast, learn different activation patterns and are more strongly characterized by localized transformations. 
As a result, the models exhibit different inductive biases and weight different aspects of the same input differently.

\citet{He_2016_CVPR} further show that ResNets of varying depth not only achieve different accuracy levels but also exhibit distinct error characteristics. 
These errors do not occur on the same inputs but are partially complementary. 
This is critical for ensemble usage, as incorrect predictions from individual models can be compensated by correct predictions from others. 
The strong performance of individual models can therefore be leveraged more effectively within an ensemble than through the deployment of a single deep model.

Since all models belong to the same architectural class and are trained under identical conditions on the same dataset, the ensemble remains homogeneous and interpretable. 
A ResNet ensemble is therefore particularly suitable as a baseline for systematically investigating ensemble effects, model diversity, and robustness.

Moreover, a ResNet ensemble can be directly compared to a single 50-layer ResNet. 
Because the 50-layer ResNet represents the strongest individual model within the ensemble, this comparison enables a clear attribution of ensemble effects. 
Performance gains achieved by the ensemble can thus be directly attributed to model combination rather than merely increased model capacity. 
This allows for a clean separation between improvements due to ensembling and those resulting solely from the use of a strong individual model.

From the perspective of ensemble learning, combining multiple models with different inductive biases is particularly well-suited for reducing variance and improving robustness. 
Shallower ResNet variants tend to rely more heavily on local and low-level visual patterns, whereas deeper ResNets learn more abstract and semantically rich representations. 
An ensemble of ResNet models with varying depths therefore enables controlled prediction diversity without altering the underlying architectural class. 
This makes ResNet-based ensembles particularly suitable as a foundation for analyzing robustness and error diversity under distribution shift.